\chapter{Melihat ke Dalam Tubuh}
\label{ch:medis}

Kuliah deep learning for medical imaging dimulai bukan dengan jaringan saraf, melainkan dengan fisika.
Naskah dasarnya membahas osilasi dan gelombang, analisis Fourier, gelombang elektromagnetik, fenomena
gelombang seperti pembiasan, interferensi, hamburan, dan penyerapan, lalu cara menghasilkan gambar
dengan sinar-X, tomografi komputer, dan kedokteran nuklir. Pilihannya masuk akal: kita tidak bisa
menilai apa yang dilakukan sebuah jaringan pada citra CT tanpa memahami dari mana angka-angka di
citra itu datang.

Citra medis berbeda dari foto biasa. Setiap piksel adalah hasil pengukuran fisik dengan satuan dan
noise yang khas, gambar sering berbentuk volume tiga dimensi, data berlabel mahal karena butuh dokter
ahli, dan kesalahan bisa berakibat fatal. Bab ini mengikuti jalan dari foton ke diagnosis.

\section{Sinar-X dan penyerapan}

Sinar-X dihasilkan dengan menembakkan elektron yang dipercepat tegangan tinggi ke sebuah anoda. Energi
foton terbesar yang bisa muncul sama dengan energi kinetik elektron, $E=h\nu=eU$, sehingga tabung dengan
tegangan 100 kilovolt menghasilkan foton sampai 100 kiloelektronvolt. Ketika berkas sinar-X melewati
tubuh, sebagian foton diserap atau dihamburkan. Penurunan intensitas di setiap irisan tipis sebanding
dengan intensitas yang datang dan dengan ketebalan irisan, $\dd I=-\mu I\dd s$, dengan $\mu$ koefisien
atenuasi jaringan. Memisahkan variabel dan mengintegrasikan memberi hukum Beer--Lambert,
\begin{equation}
I=I_0\exp\Big(-\int\mu(s)\dd s\Big).
\end{equation}
Untuk jaringan seragam dengan $\mu=0{,}2\ \mathrm{cm}^{-1}$ setebal 10 cm, hanya
$e^{-2}\approx13{,}5$ persen foton yang tembus. Tulang menyerap jauh lebih kuat dari jaringan lunak,
itulah sebabnya tulang tampak terang pada foto rontgen. Kuliah juga membahas sisi praktisnya: grid anti
hamburan yang menyaring foton yang berbelok arah, detektor digital, dan dosis radiasi, yang di Eropa
Tengah secara alami sekitar dua milisievert per tahun.

\section{Tomografi komputer}

Foto rontgen biasa menumpuk seluruh tubuh menjadi satu bayangan. Tomografi komputer mengukur banyak
proyeksi dari berbagai sudut lalu merekonstruksi peta $\mu$ di setiap titik irisan. Dari hukum
Beer--Lambert, setiap pengukuran memberi integral garis koefisien atenuasi,
\begin{equation}
p(\theta,r)=-\ln\frac{I}{I_0}=\int_{\text{garis}(\theta,r)}\mu(x,y)\dd s,
\end{equation}
dan kumpulan semua integral garis disebut transformasi Radon. Masalah utama CT adalah membalik
transformasi ini.

Cara paling naif adalah \istilah{backprojection}: setiap proyeksi ``dioleskan'' kembali sepanjang garisnya,
lalu semua sudut dijumlahkan. Hasilnya kabur, karena setiap titik terang ikut menyebar ke seluruh garis
yang melewatinya, dengan kekaburan yang meluruh seperti $1/r$. Teorema irisan pusat menjelaskan obatnya:
transformasi Fourier satu dimensi dari sebuah proyeksi sama dengan satu irisan melalui pusat dari
transformasi Fourier dua dimensi gambar. Karena irisan-irisan itu berkumpul rapat di dekat pusat dan
jarang di frekuensi tinggi, frekuensi rendah terwakili berlebihan. Mengalikan setiap proyeksi dengan
filter yang naik linear terhadap frekuensi, $|k|$, sebelum backprojection mengoreksi ketidakseimbangan
itu. Inilah \istilah{filtered backprojection} \citep{kak1988}, algoritme yang bekerja di dalam hampir semua
pemindai CT sejak pemindai pertama yang dibangun Hounsfield \citep{hounsfield1973}.

Nilai yang ditampilkan bukan $\mu$ langsung, melainkan satuan Hounsfield,
$\mathrm{HU}=1000\,(\mu-\mu_{\mathrm{air}})/\mu_{\mathrm{air}}$, dengan $\mu_{\mathrm{air}}$ koefisien
atenuasi air. Air bernilai 0 dan udara sekitar $-1000$, sedangkan tulang padat bisa beberapa ratus
sampai lebih dari seribu. Skala yang dibakukan ini membuat citra dari pemindai berbeda bisa dibandingkan,
dan menjadi alasan mengapa normalisasi intensitas untuk jaringan saraf pada CT berbeda dari normalisasi
pada foto biasa. Benda logam seperti tambalan gigi menyerap hampir semua foton dan merusak asumsi
rekonstruksi, sehingga muncul artefak bergaris yang khas.

\section{Kedokteran nuklir}

Pada skintigrafi dan PET, sumber radiasinya ada di dalam tubuh. Pasien diberi zat penanda radioaktif
yang terkumpul di jaringan dengan aktivitas tertentu, misalnya sel tumor yang memakan banyak glukosa.
Pada PET, positron yang dipancarkan bertemu elektron dan musnah menjadi dua foton yang terbang ke arah
berlawanan. Detektor yang menangkap kedua foton pada saat yang hampir bersamaan tahu bahwa peristiwanya
terjadi di sepanjang garis penghubung keduanya, sehingga rekonstruksinya kembali menjadi masalah
tomografi. Gambarnya menunjukkan fungsi, bukan anatomi, dan sering digabung dengan CT yang memberi
anatomi.

\section{Deep learning untuk citra medis}

Tugas-tugas utama di citra medis sama dengan di \cref{ch:cv}: klasifikasi, deteksi, dan segmentasi,
dengan rekonstruksi sebagai tugas tambahan yang khas. Klasifikasi lesi kulit dengan jaringan konvolusi
yang dilatih pada lebih dari seratus ribu gambar mencapai kinerja setara dokter kulit dalam percobaan
\citet{esteva2017}. Untuk segmentasi organ dan tumor, U-Net \citep{ronneberger2015} yang awalnya
dirancang untuk citra biomedis menjadi standar. nnU-Net \citep{isensee2021} menunjukkan bahwa sebagian
besar keberhasilan datang dari menyetel resep dengan benar untuk setiap himpunan data, yaitu resolusi,
normalisasi, augmentasi, dan arsitektur, dan bahwa resep itu bisa dibuat otomatis dari sifat-sifat
datanya.

Segmentasi dinilai dengan koefisien Dice, $2|A\cap B|/(|A|+|B|)$, ukuran tumpang tindih antara daerah
tebakan $A$ dan daerah yang benar $B$. Dice berkerabat dekat dengan Tanimoto di \cref{ch:hayati}: kalau
Tanimoto $T=|A\cap B|/|A\cup B|$, maka Dice $=2T/(1+T)$, karena $|A|+|B|=|A\cup B|+|A\cap B|$. Tumor yang
kecil membuat Dice peka sekali: meleset beberapa voxel saja sudah menurunkan skor dengan tajam.

Rekonstruksi juga mulai dibantu jaringan saraf. CT dosis rendah memakai lebih sedikit foton sehingga
lebih aman, tetapi gambarnya lebih berisik. Jaringan bisa membersihkan noise setelah rekonstruksi, atau
algoritme rekonstruksi iteratif bisa dibuka menjadi lapisan-lapisan yang dipelajari, seperti yang
disinggung di \cref{ch:sinyal}. Di semua penerapan ini, pengetahuan fisika tentang bagaimana citra
terbentuk tetap menjadi pegangan untuk memeriksa apakah sebuah hasil masuk akal.

\section{Batas dan tanggung jawab}

Kuliah AI dalam ilmu hayati dan kuliah tentang AI dan masyarakat sama-sama menyoroti batas sistem
seperti ini. Model bisa belajar jalan pintas, misalnya penanda di sudut gambar yang hanya muncul pada
pasien dari rumah sakit tertentu. Model yang dilatih di satu rumah sakit bisa jatuh kinerjanya di rumah
sakit lain dengan pemindai berbeda. Dan model bisa bekerja lebih buruk untuk kelompok pasien yang kurang
terwakili dalam data latih. Alat-alat XAI di \cref{ch:xai} membantu menemukan masalah seperti ini, dan
regulasi di \cref{ch:manusia} menetapkan kewajiban bagi sistem AI berisiko tinggi, termasuk alat diagnosis.

\sumberbab{Bab ini mengikuti naskah dasar Deep Learning for Medical Imaging (gelombang, analisis Fourier,
gelombang elektromagnetik, pencitraan sinar-X, tomografi komputer, kedokteran nuklir). Untuk rekonstruksi
tomografi, \citet{kak1988} adalah rujukan klasik.}
